\documentclass[12pt,a4paper]{article}

\usepackage[margin=2.5cm]{geometry}
\usepackage{booktabs}
\usepackage{tabularx}
\usepackage{array}
\usepackage{hyperref}
\usepackage{enumitem}
\usepackage{fancyhdr}
\usepackage{parskip}
\usepackage{titlesec}
\usepackage{xcolor}

\definecolor{accent}{RGB}{30,80,160}

\hypersetup{colorlinks=true, linkcolor=accent, urlcolor=accent}

\titleformat{\section}{\large\bfseries\color{accent}}{}{0em}{}[\titlerule]
\titleformat{\subsection}{\normalsize\bfseries}{}{0em}{}

\pagestyle{fancy}
\fancyhf{}
\lhead{\textbf{WanderSync} \textit{Task 1: Requirements \& Subsystems}}
\rhead{Team 15}
\rfoot{\thepage}
\renewcommand{\headrulewidth}{0.4pt}

\newcolumntype{L}[1]{>{\raggedright\arraybackslash}p{#1}}

\begin{document}

\begin{center}
  {\LARGE\bfseries Task 1: Requirements and Subsystems}\\[0.4em]
  {\large WanderSync: Dynamic Itinerary, Disruption Management and Social Travel Hub}\\[0.2em]
  {\normalsize Team 15}
\end{center}

\vspace{1em}

% ================================================================
\section{Functional Requirements}
% ================================================================

The following functional requirements define the observable behaviours the WanderSync system must
provide. Each requirement is expressed as a user story and tagged with its architectural significance.

\subsection{FR1:User Authentication and Account Management}

\textit{As a traveller, I want to register with an email and password, log in, and maintain a personal
profile, so that my itineraries and preferences are securely associated with my account.}

\textbf{Architectural significance:} Drives stateless JWT-based session management and the separation
of the Authentication Service from domain logic. Establishes security boundaries respected by every
other subsystem.

\subsection{FR2:Unified Booking Aggregation}

\textit{As a traveller, I want to import booking confirmations from multiple providers (flights, hotels,
local transport) into a single normalised view, so that all travel details are accessible in one place.}

\textbf{Architectural significance:} Requires an Adapter layer per provider with a common
\texttt{BookingRecord} interface. Adding a new provider must touch no more than 2 files outside the
adapter module:directly constraining the architecture towards low coupling and high cohesion.

\subsection{FR3:Chronological Itinerary Timeline}

\textit{As a traveller, I want to view all my bookings as a time-ordered timeline, so that I can
understand the sequence and gaps in my travel plan at a glance.}

\textbf{Architectural significance:} Requires the Itinerary Management Service to own an
authoritative, sorted view assembled from heterogeneous booking records. Timeline assembly is a
read-heavy, cacheable operation that informs the caching strategy.

\subsection{FR4:Real-Time Disruption Detection and Notification}

\textit{As a traveller, I want to be automatically notified when a disruption (flight delay,
cancellation, or severe weather) affects my itinerary, and to receive suggested alternatives, so
that I can re-plan with minimal effort.}

\textbf{Architectural significance:} This is the most architecturally significant requirement. It
mandates an Event-Driven Architecture in which the Disruption Detection Service subscribes to
external feeds via a message broker, processes events asynchronously, and publishes
\texttt{DisruptionEvent} messages consumed by the Notification and Itinerary services. This
decoupling is the primary driver of the chosen architectural style.

\subsection{FR5:Privacy-Preserving Social Travel Synchronisation}

\textit{As a traveller who opts in, I want to discover other users whose itineraries overlap with
mine in time and region, and to connect with them through a bilateral consent flow, so that I can
coordinate ride-shares or social activities while retaining control over my location data.}

\textbf{Architectural significance:} Requires geospatial indexing for efficient region-based
matching and a consent state machine (none $\to$ city-level $\to$ precise) that enforces the
privacy NFR at the data layer, not only the presentation layer.


% ================================================================
\section{Non-Functional Requirements}
% ================================================================

Each NFR is documented using the six-part quality-attribute scenario framework. The
\textit{Response Measure} column provides the verifiable acceptance criterion used during testing.

\subsection{NFR1:Performance: Response Time}

\begin{tabularx}{\textwidth}{L{3.5cm}X}
\toprule
\textbf{Scenario Element} & \textbf{Description} \\
\midrule
Source            & End user via web browser \\
Stimulus          & Navigates to the itinerary timeline page \\
Artifact          & Frontend layer and Itinerary Management Service \\
Environment       & Runtime; standard broadband connection; normal system load \\
Response          & Page renders with all booking data fully visible \\
Response Measure  & Complete render in under \textbf{3 seconds} from initial request; external API
                    calls resolved within \textbf{5 seconds} with a loading indicator shown after
                    2 seconds \\
\bottomrule
\end{tabularx}

\textbf{Architectural significance:} Drives in-process caching of itinerary timelines and
asynchronous, non-blocking API calls to external providers.

\vspace{1em}

\subsection{NFR2:Reliability: Disruption Detection Accuracy}

\begin{tabularx}{\textwidth}{L{3.5cm}X}
\toprule
\textbf{Scenario Element} & \textbf{Description} \\
\midrule
Source            & External flight-tracking or weather-alert API \\
Stimulus          & Publishes a disruption event (delay, cancellation, or weather warning) \\
Artifact          & Disruption Detection Service \\
Environment       & Runtime; system actively monitoring subscribed feeds \\
Response          & System detects the disruption, evaluates impact on active itinerary, and
                    delivers a notification with suggested alternatives \\
Response Measure  & At least \textbf{17 of 20} simulated disruption test cases correctly
                    identified and notified (\textbf{85\% accuracy}); notification delivered
                    within 10 seconds of event arrival \\
\bottomrule
\end{tabularx}

\textbf{Architectural significance:} Requires an event queue for guaranteed delivery, a retry
mechanism for transient API failures, and graceful degradation (user-facing error message, no
crash) on unrecoverable provider errors.

\vspace{1em}

\subsection{NFR3:Security: Authentication and Data in Transit}

\begin{tabularx}{\textwidth}{L{3.5cm}X}
\toprule
\textbf{Scenario Element} & \textbf{Description} \\
\midrule
Source            & Unauthenticated or malicious user \\
Stimulus          & Attempts to access a protected resource or transmit credentials \\
Artifact          & Authentication Service and API Gateway \\
Environment       & Runtime; any network condition \\
Response          & Request rejected at the gateway; user redirected to login; all
                    communication encrypted \\
Response Measure  & \textbf{100\%} of protected endpoints require a valid JWT (24-hour expiry);
                    passwords stored with \textbf{bcrypt}; all client--server traffic over
                    \textbf{HTTPS only} \\
\bottomrule
\end{tabularx}

\textbf{Architectural significance:} Drives stateless authentication at the API Gateway, centralised
token validation middleware, and TLS termination at the entry point rather than per service.

\vspace{1em}

\subsection{NFR4:Privacy: Location Disclosure Control}

\begin{tabularx}{\textwidth}{L{3.5cm}X}
\toprule
\textbf{Scenario Element} & \textbf{Description} \\
\midrule
Source            & Authenticated user browsing social matching results \\
Stimulus          & Views the profile of another user with whom no consent has been established \\
Artifact          & Social Synchronisation Service \\
Environment       & Runtime; consent state between the two users is \textit{none} \\
Response          & Only a city-level region label is returned; precise coordinates and
                    itinerary details remain withheld \\
Response Measure  & \textbf{Zero} instances of precise location leakage before bilateral opt-in,
                    verified by integration tests; location data not persisted beyond the active
                    session unless explicitly saved by the user \\
\bottomrule
\end{tabularx}

\textbf{Architectural significance:} Requires a consent state machine enforced at the data-access
layer (not merely the UI), and a geospatial abstraction that returns progressively detailed
coordinates based on consent level.

\vspace{1em}

\subsection{NFR5:Modifiability: Provider Extensibility}

\begin{tabularx}{\textwidth}{L{3.5cm}X}
\toprule
\textbf{Scenario Element} & \textbf{Description} \\
\midrule
Source            & Developer extending the system \\
Stimulus          & Integrates a new travel provider (e.g., a new airline or hotel API) \\
Artifact          & Booking Aggregation Service and its provider adapter module \\
Environment       & Development time; existing adapters and core services unchanged \\
Response          & New provider added by implementing the \texttt{BookingAdapter} interface and
                    registering it; no changes required to core services \\
Response Measure  & Changes confined to \textbf{at most 2 files} outside the dedicated adapter
                    module, verifiable by code review and diff inspection \\
\bottomrule
\end{tabularx}

\textbf{Architectural significance:} Directly mandates the Adapter design pattern with a shared
\texttt{BookingAdapter} interface and a factory or registry for dynamic provider loading.

% ================================================================
\section{Subsystem Overview}
% ================================================================

WanderSync is decomposed into seven subsystems. Each subsystem owns a well-defined responsibility
boundary and communicates with others through explicit interfaces or a message broker.

\subsection{Subsystem 1:API Gateway and Frontend Layer}

\textbf{Role:} Single external entry point for all client requests; serves the React single-page
application; enforces HTTPS and routes authenticated requests to the appropriate backend service.

\textbf{Responsibilities:}
\begin{itemize}[noitemsep]
  \item TLS termination and HTTPS enforcement for all traffic.
  \item Request routing to Authentication, Itinerary, Disruption, and Social services.
  \item JWT validation middleware applied before forwarding any protected request.
  \item Static asset delivery for the React frontend.
\end{itemize}

\subsection{Subsystem 2:Authentication and User Management Service}

\textbf{Role:} Manages user identity, credential storage, session tokens, and profile data. All
other services trust JWTs issued by this subsystem.

\textbf{Responsibilities:}
\begin{itemize}[noitemsep]
  \item User registration with bcrypt password hashing.
  \item Login, JWT issuance (24-hour expiry), and token refresh.
  \item Profile CRUD (display name, preferences, linked provider accounts).
  \item Session invalidation on explicit logout.
\end{itemize}

\subsection{Subsystem 3:Booking Aggregation Service}

\textbf{Role:} Adapts data from multiple external travel provider APIs into a canonical
\texttt{BookingRecord} schema. Shields the rest of the system from provider-specific data formats.

\textbf{Responsibilities:}
\begin{itemize}[noitemsep]
  \item Implements a \texttt{BookingAdapter} interface per provider (flight, hotel, transport).
  \item Fetches and normalises raw provider responses into \texttt{BookingRecord} objects.
  \item Forwards normalised records to the Itinerary Management Service.
  \item Enforces the extensibility NFR: a new provider requires only one new adapter file plus
        one registration entry.
\end{itemize}

\subsection{Subsystem 4:Itinerary Management Service}

\textbf{Role:} Authoritative store and assembler for a user's complete itinerary. Owns the
itinerary data store and produces the chronological timeline consumed by the frontend.

\textbf{Responsibilities:}
\begin{itemize}[noitemsep]
  \item CRUD operations on itinerary entities and their constituent booking records.
  \item Chronological timeline assembly from heterogeneous \texttt{BookingRecord} objects.
  \item Caching of assembled timelines to satisfy the 3-second load NFR.
  \item Consuming \texttt{DisruptionEvent} messages to mark affected bookings and trigger
        re-ordering.
\end{itemize}

\subsection{Subsystem 5:Disruption Detection Service}

\textbf{Role:} Continuously monitors external flight-status and weather-alert feeds. On detecting
an event that affects a stored itinerary, it evaluates the impact and publishes a
\texttt{DisruptionEvent} to the internal message broker.

\textbf{Responsibilities:}
\begin{itemize}[noitemsep]
  \item Polling or webhook subscription to flight-tracking and weather APIs.
  \item Impact evaluation: cross-referencing incoming events against active itineraries.
  \item Publishing \texttt{DisruptionEvent} messages (affected booking IDs, severity, alternatives)
        to the message broker.
  \item Maintaining an event log for the 85\% accuracy NFR to be tested against simulated cases.
\end{itemize}

\subsection{Subsystem 6:Social Synchronisation Service}

\textbf{Role:} Identifies users whose itineraries overlap in time and geographic region, manages
the bilateral consent state machine, and enforces privacy-preserving location disclosure rules.

\textbf{Responsibilities:}
\begin{itemize}[noitemsep]
  \item Geospatial indexing of anonymised itinerary waypoints for region-based matching.
  \item Consent state machine: \textit{none} $\to$ \textit{city-level} $\to$ \textit{precise},
        transitioned only by explicit bilateral opt-in actions.
  \item Returns only city-level labels to users whose consent state is \textit{none}.
  \item Purges raw location data at session end unless the user explicitly saves it.
\end{itemize}

\subsection{Subsystem 7:Notification Service}

\textbf{Role:} Dedicated consumer of \texttt{DisruptionEvent} messages from the message broker.
Delivers user-facing alerts (in-app notifications and email) to affected travellers. Decouples
the notification delivery mechanism from disruption detection and itinerary management logic,
so that new notification channels can be added without modifying either upstream service.

\textbf{Responsibilities:}
\begin{itemize}[noitemsep]
  \item Subscribing to the Redis Streams channel for \texttt{DisruptionEvent} messages as a
        dedicated consumer group.
  \item Resolving each affected user's notification preferences (in-app, email, or both) via
        the User profile store.
  \item Delivering alerts within the 10-second delivery window mandated by NFR2.
  \item Maintaining a notification delivery log for auditability and replay on missed deliveries.
\end{itemize}

% ================================================================
\section{Non-Functional Requirement 6}
% ================================================================

\subsection{NFR6:Usability: Session Continuity}

\begin{tabularx}{\textwidth}{L{3.5cm}X}
\toprule
\textbf{Scenario Element} & \textbf{Description} \\
\midrule
Source            & Authenticated traveller \\
Stimulus          & Accesses the application during a multi-leg journey whose total duration
                    exceeds 24 hours \\
Artifact          & Authentication Service and client-side session management \\
Environment       & Runtime; user mid-journey on intermittent or low-bandwidth connectivity
                    (e.g., airport Wi-Fi) \\
Response          & The 24-hour JWT expires; the application prompts re-authentication; the
                    user cannot view disruption notifications until they log in again \\
Response Measure  & Forced re-login occurs no more than once per 24-hour window; a
                    refresh-token or sliding-expiry mechanism is deferred to post-prototype
                    scope and explicitly documented as a known usability debt \\
\bottomrule
\end{tabularx}

\textbf{Architectural significance:} The stateless JWT decision (ADR-004) introduces a negative
usability trade-off for extended journeys. Mitigating this requires a refresh-token mechanism or
server-side sliding expiry, both of which re-introduce shared state and are deferred for the
prototype.

\end{document}

